\chapter{Formula Compendium (Grouped and Labeled)}\label{ch:formula-compendium}
This appendix gathers the core formulas used across Chapters 1--6.


\section*{Chapter 1 Formula Set}\label{sec:auto-056}

\paragraph{Chapter 1-F1 (Vector Serial vs Pipelined Cycles)}\label{par:auto-043}
\begin{equation}\label{eq:auto-054}
    \text{serial cycles}=nk,\qquad
    \text{pipeline cycles}\approx k+n-1.
\end{equation}


\section*{Chapter 2 Formula Set}\label{sec:auto-057}

\paragraph{Chapter 2-F1 (CPU Time Decomposition)}\label{par:auto-044}
\begin{equation}\label{eq:auto-055}
    \text{CPU time}=\text{IC}\times\text{CPI}\times\text{CT}.
\end{equation}

\paragraph{Chapter 2-F2 (Frequency Form)}\label{par:auto-045}
\begin{equation}\label{eq:auto-056}
    \text{CPU time}=\frac{\text{IC}\times\text{CPI}}{f}.
\end{equation}

\paragraph{Chapter 2-F3 (Instruction-Mix Total Cycles)}\label{par:auto-046}
\begin{equation}\label{eq:auto-057}
    \text{Total cycles}=\sum_i\left(\text{IC}_i\times\text{CPI}_i\right).
\end{equation}

\paragraph{Chapter 2-F4 (Average CPI)}\label{par:auto-047}
\begin{equation}\label{eq:auto-058}
    \text{CPI}_{avg}=\frac{\text{Total cycles}}{\text{Total instructions}}.
\end{equation}

\paragraph{Chapter 2-F5 (Speedup Definition)}\label{par:auto-048}
\begin{equation}\label{eq:auto-059}
    S=\frac{T_{old}}{T_{new}}.
\end{equation}

\paragraph{Chapter 2-F6 (Amdahl's Law)}\label{par:auto-049}
\begin{equation}\label{eq:auto-060}
    S_{Amdahl}=\frac{1}{(1-f)+\frac{f}{s}}.
\end{equation}

\paragraph{Chapter 2-F7 (Amdahl Upper Bound)}\label{par:auto-050}
\begin{equation}\label{eq:auto-061}
    S_{\text{upper}}=\frac{1}{1-f}.
\end{equation}

\paragraph{Chapter 2-F8 (Gustafson's Law)}\label{par:auto-051}
\begin{equation}\label{eq:auto-062}
    S_{Gustafson}=\alpha+(1-\alpha)N.
\end{equation}

\paragraph{Chapter 2-F9 (Arithmetic Mean for Time)}\label{par:auto-052}
\begin{equation}\label{eq:auto-063}
    \bar{T}_{arith}=\frac{1}{n}\sum_{i=1}^{n}T_i.
\end{equation}

\paragraph{Chapter 2-F10 (Harmonic Mean for Rate)}\label{par:auto-053}
\begin{equation}\label{eq:auto-064}
    \bar{R}_{harm}=\frac{n}{\sum_{i=1}^{n}\frac{1}{R_i}}.
\end{equation}


\section*{Chapter 3 Formula Set}\label{sec:auto-058}

\paragraph{Chapter 3-F1 ($(C,B,S)$ Definitions)}\label{par:auto-054}
\begin{equation}\label{eq:auto-065}
    \text{capacity}=2^C\ \text{bytes},\qquad
    \text{block size}=2^B\ \text{bytes},\qquad
    \text{associativity}=2^S\ \text{ways}.
\end{equation}

\paragraph{Chapter 3-F2 (Total Blocks and Sets)}\label{par:auto-055}
\begin{equation}\label{eq:auto-066}
    \text{blocks}=2^{C-B},\qquad
    \text{sets}=\frac{2^C}{2^B\cdot 2^S}=2^{C-B-S}.
\end{equation}
\begin{equation}\label{eq:auto-087}
    \text{lines}=\text{sets}\cdot\text{ways}=2^{C-B-S}\cdot 2^S=2^{C-B}.
\end{equation}
Direct-mapped corollary: if $S=0$, then $\text{ways}=1$ and $\text{lines}=\text{sets}$.
Set-associative corollary: if $S>0$, then $\text{lines}>\text{sets}$.

\paragraph{Chapter 3-F3 (Address-Field Widths, General $AW$)}\label{par:auto-056}
\begin{equation}\label{eq:auto-067}
    |\text{Offset}|=B,\qquad
    |\text{Index}|=C-B-S,\qquad
    |\text{Tag}|=AW-(C-S).
\end{equation}
\begin{equation}\label{eq:auto-068}
    \text{Address}=[\text{Tag}\mid\text{Index}\mid\text{Offset}].
\end{equation}

\paragraph{Chapter 3-F4 (Address-Field Widths, 32-bit Address)}\label{par:auto-057}
\begin{equation}\label{eq:auto-069}
    |\text{Offset}|=B,\qquad
    |\text{Index}|=C-B-S,\qquad
    |\text{Tag}|=32-(C-S)=32-C+S.
\end{equation}

\paragraph{Chapter 3-F5 (Way Selection Relationship)}\label{par:auto-058}
\begin{equation}\label{eq:auto-070}
    \text{ways}=2^S,\qquad
    \log_2(\text{ways})=S.
\end{equation}

\paragraph{Chapter 3-F6 (Single-Level Average Access Time)}\label{par:auto-059}
\begin{equation}\label{eq:auto-071}
    AAT=HT+MR\cdot MP.
\end{equation}

\paragraph{Chapter 3-F7 (Bit Overhead from Valid/Dirty)}\label{par:auto-060}
\begin{equation}\label{eq:auto-072}
    \text{overhead fraction}=\frac{n_v+n_d}{2^{B+3}},\qquad
    n_v=1,\quad n_d\in\{0,1\}.
\end{equation}

\paragraph{Chapter 3-F8 (Full-Block Miss Penalty Model)}\label{par:auto-061}
\begin{equation}\label{eq:auto-073}
    MP_{full}=T_1+(N_{sb}-1)T_f,\qquad N_{sb}=2^B/W.
\end{equation}

\paragraph{Chapter 3-F9 (Early-Restart Approximation)}\label{par:auto-062}
\begin{equation}\label{eq:auto-074}
    MP_{ER}\approx T_1+\frac{N_{sb}-1}{2}T_f.
\end{equation}

\paragraph{Chapter 3-F10 (Recursive Two-Level Form)}\label{par:auto-063}
\begin{equation}\label{eq:auto-075}
    AAT_{L1}=HT_{L1}+MR_{L1}\cdot AAT_{L2}.
\end{equation}

\paragraph{Chapter 3-F11 (VIPT Constraint)}\label{par:auto-064}
\begin{equation}\label{eq:auto-076}
    C-S\le P.
\end{equation}


\section*{Chapter 4 Formula Set}\label{sec:auto-059}

\paragraph{Chapter 4-F1 (Pipelined Clock Approximation)}\label{par:auto-065}
\begin{equation}\label{eq:auto-077}
    T_{clk,pipe}\approx T_s+T_{reg}.
\end{equation}

\paragraph{Chapter 4-F2 (Ideal Long-Run Pipeline Speedup)}\label{par:auto-066}
\begin{equation}\label{eq:auto-078}
    S_{ideal}\approx\frac{N T_s}{T_s+T_{reg}}.
\end{equation}

\paragraph{Chapter 4-F3 (Finite-Length Pipelined Time)}\label{par:auto-067}
\begin{equation}\label{eq:auto-079}
    T_{pipe}\approx(N+M-1)T_{clk,pipe}.
\end{equation}

\paragraph{Chapter 4-F4 (Unpipelined Time for $M$ Instructions)}\label{par:auto-068}
\begin{equation}\label{eq:auto-080}
    T_{unpiped}\approx MNT_s.
\end{equation}

\paragraph{Chapter 4-F5 (Practical CPI Model)}\label{par:auto-069}
\begin{equation}\label{eq:auto-081}
    CPI=1+\text{stalls}_{avg}.
\end{equation}

\paragraph{Chapter 4-F6 (Branch-Stall CPI Contribution)}\label{par:auto-070}
\begin{equation}\label{eq:auto-082}
    CPI\approx 1+f_{branch}\cdot \text{stall}_{branch,avg}.
\end{equation}


\section*{Chapter 5 and Simulation Formula Set}\label{sec:auto-060}

\paragraph{Chapter 5-F1 (Two-Level AAT with Explicit L1 Penalty)}\label{par:auto-071}
\begin{equation}\label{eq:auto-083}
    AAT_{L1}=HT_{L1}+MR_{L1}\cdot MP_{L1},\qquad
    MP_{L1}=AAT_{L2}.
\end{equation}

\paragraph{Chapter 5-F2 (L2 AAT Model)}\label{par:auto-072}
\begin{equation}\label{eq:auto-084}
    AAT_{L2}=HT_{L2}+MR_{L2,read}\cdot DRAM\_TIME.
\end{equation}

\paragraph{Chapter 5-F3 (DRAM Block Service Model)}\label{par:auto-073}
\begin{equation}\label{eq:auto-085}
    DRAM\_TIME=DRAM\_AT+DRAM\_AT\_PER\_WORD\cdot WORDS\_PER\_BLOCK.
\end{equation}


\section*{Chapter 6 Formula Set (ILP and Superscalar)}\label{sec:ch6-formulas}

\paragraph{Chapter 6-F1 (IPC and CPI Relationship)}\label{par:ch6-formula-1}
\begin{equation}\label{eq:ch5-formula-ipc-cpi}
    IPC=\frac{1}{CPI}.
\end{equation}

\paragraph{Chapter 6-F2 (CPI/IPC Threshold Equivalence)}\label{par:ch6-formula-2}
\begin{equation}\label{eq:ch5-formula-threshold}
    CPI<1 \iff IPC>1.
\end{equation}

\paragraph{Chapter 6-F3 (Simplified Scoreboard Fire Condition)}\label{par:ch6-formula-3}
\begin{equation}\label{eq:ch5-formula-fire}
    \text{Fire if } \neg busy(src_1)\wedge \neg busy(src_2)\wedge \neg busy(dest)\wedge \neg busy(FU_{needed}).
\end{equation}
